\documentclass[11pt,a4paper]{article}
\usepackage[T1]{fontenc}
\usepackage[utf8]{inputenc}
\usepackage{amsmath,amssymb,amsthm}
\usepackage[margin=1in]{geometry}
\usepackage{booktabs}
\usepackage[colorlinks=true,linkcolor=blue,urlcolor=blue]{hyperref}
\theoremstyle{plain}
\newtheorem{theorem}{Theorem}
\newtheorem{lemma}[theorem]{Lemma}
\newtheorem{corollary}[theorem]{Corollary}
\theoremstyle{definition}
\newtheorem{remark}[theorem]{Remark}
\newcommand{\ord}{\operatorname{ord}}
\newcommand{\Z}{\mathbb{Z}}
\title{Eventual periodicity modulo $m$ with period dividing $\varphi(m)$:\\a proof of Bala's conjecture on OEIS A002050}
\author{Adrian Perez Fontelles\\ \small Independent researcher}
\date{25 August 2026}
\begin{document}
\maketitle

\begin{abstract}
Peter Bala conjectured on OEIS A002050 that the sequence reduced modulo $k$ is eventually periodic with period dividing $\varphi(k)$. We prove it, as an instance of the general statement --- also Bala's --- that this holds for every integer sequence whose exponential generating function has the form $G(e^{x}-1)$ with $G$ an integral power series. Expanding $(e^{x}-1)^{k}$ in Stirling numbers of the second kind writes $a(n)=\sum_{k}c_{k}k!\,S(n,k)$; modulo $m$ every term with $k\ge m$ dies because $m\mid k!$, leaving a fixed finite sum, which inclusion--exclusion turns into an integer combination of the sequences $n\mapsto j^{n}$, each eventually periodic modulo $m$ with period dividing $\varphi(m)$. The identification of $G$ for this entry is carried out in Section~3.
\end{abstract}

\noindent\small 2020 Mathematics Subject Classification. 11B50, 11B73, 05A15, 11A07.\normalsize

\section{The sequence and the conjecture}

Throughout, $m$ and $k$ denote positive integers, $\varphi$ is Euler's totient function
(OEIS A000010), and $S(n,k)$ are the Stirling numbers of the second kind, with
$S(0,0)=1$ and $S(n,k)=0$ for $k>n$. We write $0^{0}=1$, and throughout
\[
y := e^{x}-1, \qquad\text{so that}\qquad e^{x}=1+y .
\]

OEIS A002050 is ``Number of simplices in barycentric subdivision of n-simplex''. It has offset $0$ and begins
\[
a(0),\dots,a(7)\;=\;0, 1, 5, 25, 149, 1081, 9365, 94585,\ \dots
\]
and it carries the following comment:
\begin{quote}\small
\textbf{Conjecture:} Let k be a positive integer. The sequence obtained by reducing a(n) modulo k is eventually periodic with the period dividing phi(k) = A000010(k). For example, modulo 9 we obtain the sequence [0, 1, 5, 7, 5, 1, 5, 4, 5, 7, 5, 1, 5, 4, 5, 7, 5, 1, 5, 4, 5, 7, 5, ...], with an apparent period of 6 = phi(9) beginning at a(5). - \emph{Peter Bala}, Aug 03 2023
\end{quote}
As of the ``Last modified'' line on the live entry (2026-08-05, revision 97) the
statement is still recorded as a conjecture, and no proof appears among the entry's links,
formulas or programs.

Bala also states the conjecture in a general form: that the property should hold for
every integer sequence whose exponential generating function has the shape $G(e^{x}-1)$
with $G$ an integral power series. We prove that general statement as
Theorem~\ref{thm:main} below, and then obtain the present entry as
Corollary~\ref{cor:this} by identifying its $G$ in Section~\ref{sec:apply}.

\section{The theorem}

\begin{theorem}\label{thm:main}
Let $G(t)=\sum_{k\ge0}c_{k}t^{k}$ be a power series with integer coefficients, and define
integers $a(n)$ by $\sum_{n\ge0}a(n)x^{n}/n!=G(e^{x}-1)$. Then for every $m\ge1$ the
sequence $\bigl(a(n)\bmod m\bigr)_{n\ge0}$ is eventually periodic, with period dividing
$\varphi(m)$.
\end{theorem}

The proof uses three standard facts.

\begin{lemma}\label{lem:stirling}
$(e^{x}-1)^{k}=k!\sum_{n\ge k}S(n,k)x^{n}/n!$, and hence
$a(n)=\sum_{k=0}^{n}c_{k}\,k!\,S(n,k)$ for every $n\ge0$. In particular every $a(n)$ is an
integer.
\end{lemma}

\begin{proof}
The first identity is the exponential generating function of the Stirling numbers of the
second kind. Multiplying by $c_{k}$ and summing over $k$ is legitimate in $\Z[[x]]$
because $(e^{x}-1)^{k}$ has order $k$ in $x$, so each coefficient of $x^{n}$ receives
contributions from only finitely many $k$; comparing coefficients of $x^{n}/n!$ gives the
formula, and the sum terminates because $S(n,k)=0$ for $k>n$.
\end{proof}

\begin{lemma}\label{lem:power}
$k!\,S(n,k)=\sum_{j=0}^{k}(-1)^{k-j}\binom{k}{j}j^{\,n}$ for all $n,k\ge0$.
\end{lemma}

\begin{proof}
Both sides count the surjections from an $n$-set onto a $k$-set: the left directly, the
right by inclusion--exclusion over the points of the codomain that are missed. For $n=0$
both sides equal $[k=0]$, using $0^{0}=1$.
\end{proof}

\begin{lemma}\label{lem:jn}
For integers $m\ge1$ and $j\ge0$, the sequence $\bigl(j^{\,n}\bmod m\bigr)_{n\ge0}$ is
eventually periodic with period dividing $\varphi(m)$.
\end{lemma}

\begin{proof}
Write $m=m_{1}m_{2}$, where $m_{1}$ is the product of those prime powers $p^{v_{p}(m)}$ of
$m$ whose prime $p$ divides $j$, and $m_{2}=m/m_{1}$. Then $\gcd(m_{1},m_{2})=1$ and
$\gcd(j,m_{2})=1$.

Every prime dividing $m_{1}$ divides $j$, so $m_{1}\mid j^{\,n}$ once $n\ge\log_{2}m$;
thus $j^{\,n}\equiv0 \pmod{m_{1}}$ for all large $n$, a sequence of eventual period $1$.
Modulo $m_{2}$ the integer $j$ is a unit, so $j^{\,n}$ is purely periodic with period
$\ord_{m_{2}}(j)$, a divisor of $\varphi(m_{2})$. By the Chinese remainder theorem
$j^{\,n} \bmod m$ is eventually periodic with period dividing
$\operatorname{lcm}(1,\ord_{m_{2}}(j))=\ord_{m_{2}}(j)$. Finally $m_{2}\mid m$ gives
$\varphi(m_{2})\mid\varphi(m)$, so that period divides $\varphi(m)$.
\end{proof}

\begin{proof}[Proof of Theorem~\ref{thm:main}]
Fix $m\ge1$. Because $m\mid m!$ we have $m\mid k!$ for every $k\ge m$, so in
Lemma~\ref{lem:stirling} every term with $k\ge m$ is divisible by $m$. Hence for $n\ge m$
\begin{equation}\label{eq:trunc}
a(n)\;\equiv\;\sum_{k=0}^{m-1}c_{k}\,k!\,S(n,k) \pmod m ,
\end{equation}
a sum whose number of terms no longer depends on $n$. This is the only place the
hypothesis $G\in\Z[[t]]$ is used: it keeps the $c_{k}$ integral, so no denominator can
cancel the factor $k!$ that supplies the divisibility.

Substituting Lemma~\ref{lem:power} and exchanging the two finite sums,
\begin{equation}\label{eq:expo}
a(n)\;\equiv\;\sum_{j=0}^{m-1}w_{j}\,j^{\,n} \pmod m ,
\qquad
w_{j}:=\sum_{k=j}^{m-1}(-1)^{k-j}\binom{k}{j}c_{k}\in\Z ,
\end{equation}
for all $n\ge m$; the weights $w_{j}$ depend on $m$ and $G$ but not on $n$. By
Lemma~\ref{lem:jn} each sequence $n\mapsto j^{\,n}\bmod m$ is eventually periodic with
period dividing $\varphi(m)$. A finite $\Z$-linear combination of eventually periodic
sequences is eventually periodic with period dividing the least common multiple of their
periods, and the least common multiple of finitely many divisors of $\varphi(m)$ again
divides $\varphi(m)$. So the right-hand side of \eqref{eq:expo}, and therefore $a(n)$, is
eventually periodic modulo $m$ with period dividing $\varphi(m)$.
\end{proof}

\begin{remark}
``Eventually'' cannot be dropped, and the divisor need not be $\varphi(m)$ itself.
Equation \eqref{eq:trunc} is asserted only for $n\ge m$, and Lemma~\ref{lem:jn} adds a
further preperiod coming from the primes shared by $j$ and $m$; meanwhile the exact period
is the least common multiple of the orders $\ord_{m_{2}}(j)$ over those $j$ whose weight
$w_{j}$ survives modulo $m$, which is often a proper divisor of $\varphi(m)$. Both effects
are visible in the table of Section~\ref{sec:ver}.
\end{remark}

\section{Application to A002050}\label{sec:apply}

\begin{corollary}\label{cor:this}
The conjecture quoted in Section~1 is true: for every $m\ge1$ the sequence
A002050$(n) \bmod m$ is eventually periodic with period dividing $\varphi(m)$.
\end{corollary}

\begin{proof}
By Theorem~\ref{thm:main} it suffices to exhibit the entry's exponential generating
function in the form $G(e^{x}-1)$ with $G$ integral. The posted e.g.f.\ is $(e^{2x}-e^{x})/(2-e^{x})$. With $e^{x}=1+y$ the numerator is $(1+y)^{2}-(1+y)=y(1+y)$ and the denominator is $1-y$. Therefore
\[
\sum_{n\ge0}a(n)\frac{x^{n}}{n!}=G(y),\qquad G(t)=\frac{y(1+y)}{1-y} ,
\]
whose coefficients are a polynomial over $1-y$, all integers. So $G\in\Z[[t]]$ and
Theorem~\ref{thm:main} applies.
\end{proof}

Concretely, the first coefficients of $G$ are
\[
c_{0},c_{1},c_{2},\dots \;=\; 0, 1, 2, 2, 2, 2, 2, \dots ,
\]
and these are exactly what the inversion of Section~\ref{sec:ver} recovers from the
entry's published terms.

\section{Verification}\label{sec:ver}

Every number quoted here is an output of a script written separately from this note.

First, the identification of $G$ was checked against the entry rather than assumed. If
$\sum_{n}a(n)x^{n}/n!=G(e^{x}-1)$ then substituting $x=\log(1+y)$ gives
\[
c_{k}=\frac{1}{k!}\sum_{n}s(k,n)\,a(n),
\]
with $s$ the signed Stirling numbers of the first kind. Applying this to the
$19$ terms published on A002050 returns integers at every $k$, and the values
agree with the coefficients of the $G$ named in Section~\ref{sec:apply}.

Second, the conclusion itself was tested. From the recovered $c_{k}$ we generated
$a(n)\bmod m$ for $n\le400$ using $a(n)=\sum_{k}c_{k}k!\,S(n,k)$ with an integer
recurrence for $S(n,k)$ --- a computation that touches neither the generating function nor
the published terms --- then measured the exact eventual period over the last $160$
values. The results:

\begin{center}
\begin{tabular}{@{}lcccccccccc@{}}
\toprule
$m$  & $5$ & $7$ & $8$ & $9$ & $11$ & $12$ & $13$ & $16$ & $25$ & $27$\\
$\varphi(m)$  & $4$ & $6$ & $4$ & $6$ & $10$ & $4$ & $12$ & $8$ & $20$ & $18$\\
exact period  & $4$ & $6$ & $2$ & $6$ & $10$ & $2$ & $12$ & $2$ & $20$ & $18$\\
\bottomrule
\end{tabular}
\end{center}

\noindent
In every column the exact period divides $\varphi(m)$, as Corollary~\ref{cor:this}
asserts.

\section{Relation to earlier work}

A related note of Bala (\emph{Integer sequences that become periodic on reduction modulo
$k$ for all $k$}, December 2017, linked from OEIS A047974) proves a different theorem, for
generating functions of the shape $F(x)\exp(xG(x))$ and with the stronger conclusion
$a(n+k)\equiv a(n)\pmod{k}$, i.e.\ pure periodicity with period dividing $k$. Neither its
hypothesis nor its conclusion covers the statement proved here, and the argument of
Section~2 is independent of it.

\begin{thebibliography}{9}
\bibitem{oeis} The OEIS Foundation, \emph{The On-Line Encyclopedia of Integer Sequences},
\url{https://oeis.org}, 2026. Sequence A002050; conjecture of P.\ Bala.
\bibitem{bala} P.\ Bala, \emph{Integer sequences that become periodic on reduction modulo
$k$ for all $k$}, December 2017; linked from OEIS A047974.
\bibitem{comtet} L.\ Comtet, \emph{Advanced Combinatorics}, Reidel, 1974, Chapter V.
\end{thebibliography}
\end{document}
